% TOP_PROBLEM: {"problemId": "problem.schanuel-conjecture", "canonicalTitle": "Schanuel's conjecture", "releaseRank": 28, "primaryDomain": "number_theory_arithmetic_geometry", "primaryDomainLabel": "Number theory & arithmetic geometry", "displayStatus": "open", "releaseStatus": "open"}
% !TeX program = lualatex
% ATTEMPT_STATUS: unresolved
% Research continuation by GPT_6_Astra_Ultra, 27 September 2026.
% Catalogue statements and existing sources preserved verbatim.
\documentclass[11pt]{article}
\usepackage[margin=25mm]{geometry}
\usepackage{fontspec}
\setmainfont{Latin Modern Roman}
\usepackage{amsmath,amssymb,mathtools}
\usepackage{unicode-math}
\setmathfont{Latin Modern Math}
\usepackage{xurl,hyperref}
\hypersetup{colorlinks=true,urlcolor=blue}
\setcounter{secnumdepth}{0}
\setlength{\parindent}{0pt}
\setlength{\parskip}{5pt}
\setlength{\emergencystretch}{3em}
\begin{document}
\section*{\texorpdfstring{Schanuel's conjecture}{Schanuel's conjecture}}\label{28}
\subsection{Definitions and mathematical statement}
Numbers $z_1,\ldots,z_n$ are ${\mathbb{Q}}$-linearly independent when $\sum_jq_jz_j=0$ with $q_j\in{\mathbb{Q}}$ forces all $q_j=0$. The transcendence degree of a field extension is the maximum size of an algebraically independent subset. Schanuel's assertion is
\[
 \dim_{{\mathbb{Q}}}\operatorname{span}_{{\mathbb{Q}}}\{z_1,\ldots,z_n\}=n
 \quad\Longrightarrow\quad
 \operatorname{trdeg}_{{\mathbb{Q}}}{\mathbb{Q}}(z_1,\ldots,z_n,e^{z_1},\ldots,e^{z_n})\ge n.
\]
It is quantified over every $n\ge1$ and every such complex tuple; the $z_j$ themselves need not be algebraic.
\par\smallskip\textit{Formulation and concept references:} \href{https://www.mdpi.com/2227-7390/9/7/717}{[S1]}.

\subsection{Short English statement}
Do linearly independent complex numbers, together with their exponentials, always contain at least as much algebraic independence as Schanuel predicts?
\subsection{Research attempt}
\paragraph{Status and scope.}
This is an unresolved attempt, dated 27 September 2026. The aim is to
identify exactly what a reduction to independent inputs accomplishes, prove
the algebraic-input case using a named classical theorem, and test a proposed
induction on exponential towers. The elementary reductions below are not
claimed to be new. None gives the missing transcendence estimate for arbitrary
complex inputs.

The external theorem used is Lindemann--Weierstrass: exponentials of distinct
algebraic numbers are linearly independent over the algebraic numbers. A
primary proof is [E1]. Ax's theorem [E2] concerns exponentiation in differential
fields. Its functional scope is different from the numerical assertion here;
no specialization principle is assumed. The power-tower reference [S1] is
retained as context, not used to turn a conditional conclusion into a theorem.

\paragraph{1. Removing dependent coordinates without losing exponentials.}
For a finite tuple $z=(z_1,\ldots,z_m)$, put
\[
 K(z)=\mathbb Q(z_1,\ldots,z_m,e^{z_1},\ldots,e^{z_m}),\qquad
 \delta(z)=\operatorname{trdeg}_{\mathbb Q}K(z)
       -\dim_{\mathbb Q}\langle z_1,\ldots,z_m\rangle.
\]
Choose a basis $b_1,\ldots,b_r$ of this rational span from among the $z_i$.
For a common positive denominator $D$ there are integers $a_{ij}$ with
$Dz_i=\sum_j a_{ij}b_j$. Hence
\[
 (e^{z_i})^D=\prod_{j=1}^r(e^{b_j})^{a_{ij}}.
 \tag{28.1}
\]
The right side belongs to $K(b)$, including when some exponents are negative,
because exponentials never vanish. Each $z_i$ already belongs to $K(b)$ and
each $e^{z_i}$ is algebraic over it. Conversely $K(b)\subseteq K(z)$ because
the basis was selected from the tuple. Thus $K(z)/K(b)$ is algebraic and
\[
 \operatorname{trdeg}_{\mathbb Q}K(z)
 =\operatorname{trdeg}_{\mathbb Q}K(b),\qquad \delta(z)=\delta(b).
 \tag{28.2}
\]
There is no logarithm branch choice in this proof: (28.1) follows from the
single-valued exponential addition law. The assertion for independent tuples
is therefore equivalent to $\delta(z)\geq0$ for every finite tuple, including
dependent ones. Rank zero consists of zero inputs and gives $K(z)=\mathbb Q$.

This reformulation is useful bookkeeping, but it supplies no lower bound on
the transcendence degree. In particular, selecting a rational basis does not
make its entries algebraic. Any argument applying Lindemann--Weierstrass
immediately after this selection would be using an unproved hypothesis.

\paragraph{2. Two unconditional subcases.}
For one nonzero complex number $z$, the conjecture follows at once. If $z$
is transcendental, the field already contains one transcendental element.
If $z$ is algebraic, Lindemann--Weierstrass applied to $0,z$ makes $e^z$
transcendental. This establishes the entire rank-one case, not just real or
algebraic rank-one inputs.

Now suppose $b_1,\ldots,b_r$ are algebraic and rationally independent. Let
$P\in\overline{\mathbb Q}[X_1,\ldots,X_r]$ be a nonzero polynomial, written
with its distinct monomials as
\[
 P(X)=\sum_{\nu\in A}c_\nu X_1^{\nu_1}\cdots X_r^{\nu_r},
 \qquad c_\nu\ne0.
\]
At the exponentials its value is
\[
 P(e^{b_1},\ldots,e^{b_r})
 =\sum_{\nu\in A}c_\nu\exp(\nu_1b_1+\cdots+\nu_rb_r).
 \tag{28.3}
\]
The exponents are algebraic and distinct: equality for two indices would give
the nontrivial rational relation $(\nu-\nu')\cdot b=0$.
Lindemann--Weierstrass therefore makes (28.3) nonzero. Consequently
$e^{b_1},\ldots,e^{b_r}$ are algebraically independent, and
$\operatorname{trdeg}_{\mathbb Q}K(b)=r$ since the inputs are algebraic.
Together with (28.2), this proves the rank formulation for every tuple of
algebraic numbers. This familiar corollary is a genuine subcase, with its
dependence on the classical theorem explicit.

The exact point of failure for transcendental inputs is visible in (28.3):
its exponents need no longer be algebraic. Their distinctness survives, but
distinct complex exponents alone give no nonvanishing theorem for their
exponential values. For example $0$ and $2\pi i$ are distinct and both
exponentials equal $1$. That example is a stress test of the proposed use of
the theorem, not a counterexample to Schanuel: the input $2\pi i$ itself
contributes transcendence.

\paragraph{3. What a smallest counterexample would look like.}
Suppose, conditionally, that the conjecture fails, and let $r$ be its smallest
failing rank. The rank-one argument gives $r\geq2$. Choose an independent
counterexample $b_1,\ldots,b_r$. Each $(r-1)$-coordinate subtuple satisfies
the conjecture by minimality, so its field has transcendence degree at least
$r-1$. Its field is a subfield of $K(b)$, whereas failure says the latter
degree is less than $r$. Therefore necessarily
\[
 \operatorname{trdeg}_{\mathbb Q}K(b)=r-1,\qquad \delta(b)=-1.
 \tag{28.4}
\]
More generally every rational subspace of the span of dimension less than
$r$ has nonnegative predimension, by choosing a basis and using minimality.
Thus a least-rank obstruction would be a deficit of exactly one, invisible
on all lower-dimensional rational subspaces. This narrows a hypothetical
obstruction but does not rule it out. In particular, transcendence degree
need not increase when a new pair $(b_r,e^{b_r})$ is adjoined; proving that it
must increase is precisely where an induction would require new input.

One can also see why merely counting $2r$ generators is ineffective. A field
generated by many named numbers can have much smaller transcendence degree
because of polynomial relations. Rational independence excludes linear
relations between the inputs, not arbitrary relations mixing inputs and
outputs. Equation (28.4) does not supply an algorithm for recognizing that
mixing from decimal approximations.

\paragraph{4. A tower induction, with its hypothesis exposed.}
Define real numbers $a_0=1$ and $a_{j+1}=e^{a_j}$. Assume Schanuel's
conjecture. Then $a_1,\ldots,a_m$ are algebraically independent for every
$m\geq1$. Here is the full induction. For $m=1$, $a_1=e$ is transcendental
by the preceding theorem. Suppose $a_1,\ldots,a_{m-1}$ are algebraically
independent. Any rational relation among $1,a_1,\ldots,a_{m-1}$ is a
polynomial relation in the latter $m-1$ numbers, so all its coefficients
vanish. Apply Schanuel to that $m$-tuple. The resulting field is exactly
\[
 \mathbb Q(1,a_1,\ldots,a_{m-1},
                 e,e^{a_1},\ldots,e^{a_{m-1}})
 =\mathbb Q(a_1,\ldots,a_m).
\]
It has transcendence degree at least $m$, and it is generated by $m$ numbers,
so those generators are algebraically independent.

This is a conditional consequence, not an induction proving Schanuel.
Already the step from $e$ to $(e,e^e)$ invokes the conjecture on $(1,e)$,
whose second input is transcendental. Positivity and rapid growth of the
tower do not justify that invocation. The same distinction appears in the
tuple $(1,i\pi)$: it is rationally independent by real and imaginary parts,
and Schanuel would give transcendence degree at least two for
$\mathbb Q(i\pi,e,-1)$. After adjoining the algebraic number $i$, this is
equivalent to algebraic independence of $e$ and $\pi$. Their separate
transcendence does not imply that conclusion.

\paragraph{5. A functional argument cannot be specialized for free.}
It is possible to prove directly that the functions $t,e^t$ are algebraically
independent over $\mathbb C$. Suppose
$\sum_{j=0}^M P_j(t)e^{jt}=0$ with polynomial coefficients and $P_M\ne0$.
Divide by $e^{Mt}$ and take real $t\to+\infty$. Every term with $j<M$
tends to zero, forcing the nonzero polynomial $P_M(t)$ to tend to zero,
which is impossible. Thus the meromorphic function field
$\mathbb C(t,e^t)$ has transcendence degree two over $\mathbb C$.
In particular, $\mathbb Q(t,e^t)$ has transcendence degree two over
$\mathbb Q$ as well, so the next comparison can keep the base field fixed.

At $t=1$, however, the corresponding numerical field $\mathbb Q(1,e)$ has
transcendence degree only one over $\mathbb Q$. At $t=0$ both values are
algebraic. These observations do not contradict the numerical conjecture;
they show that independence of functions is not a license to retain the
same degree at a prescribed point. Ax's differential theorem [E2] should
therefore not be inserted at the missing numerical induction step. A valid
specialization theorem would need additional arithmetic hypotheses and a
proof that the chosen point satisfies them.

\paragraph{Verification and remaining gap.}
The companion exact-arithmetic checks use the coefficient vectors of
$\sqrt2,\sqrt3,(\sqrt2+\sqrt3)/2$. They confirm rational rank two and the
cleared relation $2z_3-z_1-z_2=0$, which gives
$(e^{z_3})^2=e^{z_1}e^{z_2}$. The independence of $\sqrt2,\sqrt3$ used to
interpret these vectors is elementary: a nontrivial rational dependence
would make $\sqrt{3/2}$ rational, contradicting the parity of prime
valuations in a rational square. The script also checks distinct exponent
vectors for all pairs of monomial indices between zero and eight in each
coordinate. These finite checks test denominator handling and monomial
bookkeeping only; the general injectivity proof is the argument in (28.3).
No finite search or numerical precision is reported as a transcendence test.

The outcome is an exact rank reduction, the full rank-one and algebraic-input
subcases, a description of any least-rank failure, and explicit conditional
consequences. The first unresolved step is a lower bound for
$\operatorname{trdeg}\mathbb Q(b,e^b)$ when the independent inputs include
transcendental numbers and their known transcendence does not already reach
their rational rank. Neither the basis reduction, the tower construction,
nor differential algebra fills that step. Schanuel's conjecture remains
\textbf{UNRESOLVED}.

\subsection{Sources}
\begin{itemize}
\item[S1]Schanuel's Conjecture and the Transcendence of Power Towers.
\url{https://www.mdpi.com/2227-7390/9/7/717}.
\textit{Formulation source.}
\item[S2]Combining the conjectures of Schanuel and Zilber-Pink.
\url{https://content.ems.press/assets/public/full-texts/serials/rlm/36/3/14299588/online/10.4171-rlm-1086.pdf}.
\textit{Further reference; not a proof endorsement.}
\item[S3]Schanuel Property for Elliptic and Quasi-Elliptic Functions.
\url{https://webusers.imj-prg.fr/~michel.waldschmidt/articles/pdf/SchanuelPropertyAlmostAll.pdf}.
\textit{Further reference; not a proof endorsement.}

\item[E1]F. Beukers, J. P. B\'ezivin, and P. Robba,
\emph{An Alternative Proof of the Lindemann--Weierstrass Theorem},
American Mathematical Monthly 97 (1990), no. 3, 193--197.
\url{https://doi.org/10.2307/2324683}.
\textit{Primary proof of the classical nonvanishing theorem used in the
algebraic-input case. Full-text scan inspected at
\url{https://datuan5pdes.wordpress.com/wp-content/uploads/2016/12/10-23072324683.pdf}.
Consulted 27 September 2026.}
\item[E2]James Ax, \emph{On Schanuel's conjectures},
Annals of Mathematics 93 (1971), no. 2, 252--268.
\url{https://annals.math.princeton.edu/1971/93-2/p03}.
\textit{Primary research reference for the differential-algebraic setting;
the publisher's abstract was inspected to check scope. No numerical
specialization theorem is inferred. Consulted 27 September 2026.}
\end{itemize}
\end{document}
